\chapter[2018 --- Allison et al.]{2018: James P. Allison, Tasuku Honjo}
\label{ch:2018_Allison}

\section*{Main Contribution}

\textit{Discovered cancer therapy by inhibiting negative immune regulation (checkpoint therapy), unleashing the immune system to attack cancer cells---a new pillar of cancer treatment.}\cite{nobel-2018,lecture-2018}

\section{Official Citation}

for their discovery of cancer therapy by inhibition of negative immune regulation

\section{Background and Historical Context}

The 2018 Nobel Prize in Physiology or Medicine recognized James P. Allison and Tasuku Honjo ``for their discovery of cancer therapy by inhibition of negative immune regulation.'' Their work established the field of immune checkpoint therapy---a significant approach to cancer treatment that harnesses the body's own immune system to recognize and destroy cancer cells. Checkpoint inhibitors have transformed the treatment of many types of cancer and have given hope to patients with previously untreatable malignancies.

The idea that the immune system might be able to recognize and destroy cancer cells---known as cancer immunosurveillance---was first proposed by Paul Ehrlich in the early twentieth century and was developed into a more formal theory by Frank Macfarlane Burnet and Lewis Thomas in the 1950s and 1960s. According to this theory, the immune system constantly monitors the body for the presence of abnormal cells, including cancer cells, and eliminates them before they can develop into clinically apparent tumors. The concept was supported by observations that immunosuppressed individuals had higher rates of cancer, and that tumors often expressed abnormal proteins (tumor-associated antigens) that could be recognized by the immune system.

However, despite the theoretical appeal of cancer immunosurveillance, attempts to develop effective immunotherapies for cancer had been largely unsuccessful for decades. Early approaches, including cancer vaccines and adoptive T cell therapy, showed limited clinical efficacy, and many oncologists dismissed the idea that the immune system could be harnessed to treat established cancers. The prevailing view was that the immune system was either unable to recognize cancer cells or was actively suppressed by the tumor. This pessimism was not entirely unfounded---cancers had evolved numerous mechanisms to evade immune detection, including downregulating MHC molecules, secreting immunosuppressive factors, recruiting regulatory immune cells, and exploiting the very inhibitory pathways that normally keep the immune system in check.

The key insight that led to the development of immune checkpoint therapy was the recognition that the immune system possesses natural ``brakes'' or ``checkpoints'' that prevent excessive immune activation and autoimmunity. These checkpoints are essential for maintaining immune homeostasis and preventing the immune system from attacking the body's own tissues. However, they can also be exploited by cancer cells, which activate these inhibitory pathways to evade immune destruction. The therapeutic strategy of checkpoint blockade---releasing these brakes to unleash the immune system against cancer---was the conceptual breakthrough that led to the development of a new class of cancer immunotherapies. This concept was supported by observations that tumors expressing high levels of checkpoint ligands (such as PD-L1) tended to have a worse prognosis, suggesting that checkpoint pathways were actively protecting tumors from immune attack.

\section{The Discovery/Contribution}

\textbf{James P. Allison:} Allison, an American immunologist born in 1948 in Alice, Texas, made his foundational discovery while working at the University of California, Berkeley, and later at the Memorial Sloan-Kettering Cancer Center. In the late 1980s, Allison was studying the T cell receptor---the molecular complex on the surface of T cells that recognizes antigens presented by major histocompatibility complex (MHC) molecules. His work on the structure and function of the T cell receptor led him to investigate the role of a surface protein called CTLA-4 (cytotoxic T-lymphocyte-associated protein 4) in T cell regulation.

CTLA-4 was known to be a negative regulator of T cell activation---it acted as an inhibitory receptor that dampened T cell responses after they were initiated. Allison hypothesized that if CTLA-4 was blocking the anti-tumor activity of T cells, then blocking CTLA-4 with an antibody could unleash the immune system against cancer. In a landmark experiment published in 1996, Allison demonstrated that an antibody against CTLA-4 (anti-CTLA-4) could inhibit the growth of transplanted tumors in mice by enhancing the anti-tumor activity of T cells.

Allison's anti-CTLA-4 antibody, which he developed while at Berkeley, became the prototype for the checkpoint inhibitor drug ipilimumab (Yervoy), which was approved by the U.S. FDA in 2011 for the treatment of advanced melanoma. Ipilimumab was the first checkpoint inhibitor to receive regulatory approval and represented a new paradigm in cancer treatment. Clinical trials demonstrated that ipilimumab could produce durable responses in patients with advanced melanoma, a cancer that had previously been largely resistant to treatment, with some patients achieving long-term survival. This was particularly remarkable because ipilimumab appeared to produce a fundamentally different type of response compared to conventional chemotherapy---rather than causing rapid tumor shrinkage followed by relapse, checkpoint inhibitors could produce delayed but durable responses, with some patients surviving for more than a decade after treatment. This ``tail of the survival curve'' pattern became a hallmark of checkpoint inhibitor therapy and suggested that the immune system, once unleashed, could provide ongoing surveillance against cancer recurrence.

The mechanism of action of ipilimumab involves blocking the interaction between CTLA-4 on T cells and its ligands B7-1 (CD80) and B7-2 (CD86) on antigen-presenting cells. Under normal conditions, CTLA-4 competes with the co-stimulatory receptor CD28 for binding to B7-1 and B7-2, and its higher affinity for these ligands allows it to dominate the regulatory balance, dampening T cell activation. By blocking CTLA-4, ipilimumab tips the balance in favor of CD28-mediated co-stimulation, enhancing T cell activation and proliferation, and promoting anti-tumor immune responses.

\textbf{Tasuku Honjo:} Honjo, a Japanese immunologist born in 1942 in Kyoto, Japan, made his foundational discovery while working at Kyoto University. In 1992, Honjo and his colleagues identified a novel protein that was expressed on the surface of activated T cells and that appeared to play a role in immune regulation. They named this protein PD-1 (programmed cell death protein 1), based on its initial characterization as a molecule involved in programmed cell death.

Subsequent research by Honjo and others revealed that PD-1 was actually an inhibitory receptor---a negative regulator of T cell function that played a role in limiting immune responses in peripheral tissues. PD-1 was shown to interact with two ligands, PD-L1 (programmed cell death ligand 1) and PD-L2 (programmed cell death ligand 2), which were expressed on a wide range of cells, including tumor cells. The engagement of PD-1 by PD-L1 or PD-L2 on tumor cells delivered an inhibitory signal to T cells, effectively shutting down their anti-tumor activity.

Honjo's discovery of the PD-1 pathway provided the basis for the development of a second class of checkpoint inhibitor drugs---anti-PD-1 and anti-PD-L1 antibodies. Nivolumab (Opdivo), an anti-PD-1 antibody developed by Bristol-Myers Squibb, and pembrolizumab (Keytruda), another anti-PD-1 antibody developed by Merck, were approved by the U.S. FDA in 2014 for the treatment of advanced melanoma. These drugs have since been approved for the treatment of numerous other cancer types, including non-small cell lung cancer, renal cell carcinoma, head and neck cancer, bladder cancer, and Hodgkin lymphoma.

The mechanism of action of anti-PD-1 antibodies involves blocking the interaction between PD-1 on T cells and PD-L1 on tumor cells, thereby preventing the tumor from shutting down T cell activity. By releasing this brake on T cell function, anti-PD-1 antibodies enable the immune system to recognize and destroy cancer cells. The efficacy of anti-PD-1 therapy has been particularly striking in cancers with high mutational burdens, such as melanoma and lung cancer, where the abundance of tumor-specific mutations generates a large number of neoantigens that can be recognized by T cells.

\section{Impact on Modern Medicine}

The impact of Allison and Honjo's discoveries on modern medicine has been transformative. Checkpoint inhibitors have revolutionized the treatment of cancer and have fundamentally changed the way oncologists think about and treat the disease.

Since the approval of ipilimumab in 2011, checkpoint inhibitors have been approved for the treatment of more than 20 different types of cancer. The list of approved indications continues to grow as clinical trials demonstrate efficacy in additional cancer types. Checkpoint inhibitors are now a standard component of cancer treatment for many malignancies and have significantly improved survival outcomes for patients with advanced cancer.

The clinical impact of checkpoint inhibitors has been particularly dramatic in certain cancer types. In advanced melanoma, which had a five-year survival rate of less than 5\% before the advent of immunotherapy, checkpoint inhibitors have achieved five-year survival rates of 30--40\% in some patient populations. In non-small cell lung cancer, checkpoint inhibitors have improved survival compared to chemotherapy in both the first-line and second-line settings. In Hodgkin lymphoma, anti-PD-1 antibodies have achieved response rates exceeding 60\% in patients who had failed multiple prior therapies.

The combination of checkpoint inhibitors targeting different pathways has shown even greater efficacy than single-agent therapy. The combination of ipilimumab (anti-CTLA-4) and nivolumab (anti-PD-1) has been shown to produce superior outcomes compared to either agent alone in several cancer types, including melanoma and renal cell carcinoma. However, combination therapy is also associated with increased immune-related adverse events, requiring careful management.

Checkpoint inhibitors have also changed the way cancer is classified and staged. The concept of ``tumor-agnostic'' therapy---treating cancers based on molecular characteristics rather than tissue of origin---has been validated by the approval of pembrolizumab for tumors with high microsatellite instability (MSI-H) or mismatch repair deficiency (dMMR), regardless of the primary tumor site. This represents a paradigm shift in oncology, moving away from a tissue-based approach to a molecular-based approach to cancer treatment.

The development of biomarkers for predicting response to checkpoint inhibitors is an active area of research. Tumor mutational burden (TMB), PD-L1 expression, microsatellite instability status, and immune gene expression signatures are being evaluated as potential predictive biomarkers. The identification of reliable biomarkers would enable more precise selection of patients who are most likely to benefit from checkpoint inhibitor therapy, optimizing treatment outcomes and reducing unnecessary toxicity.

James P. Allison and Tasuku Honjo were awarded the Nobel Prize in Physiology or Medicine in 2018 ``for their discovery of cancer therapy by inhibition of negative immune regulation.'' Their work has saved millions of lives and has fundamentally changed the landscape of cancer treatment.

\section{Legacy}

The legacy of Allison and Honjo's work extends far beyond the immediate clinical applications of checkpoint inhibitors. Their discoveries have established the field of cancer immunotherapy as a major pillar of cancer treatment, alongside surgery, chemotherapy, and radiation therapy. The success of checkpoint inhibitors has validated the concept of cancer immunosurveillance and has inspired a renaissance in cancer immunology research.

Allison's discovery of CTLA-4 as an immune checkpoint and Honjo's discovery of the PD-1 pathway have revealed fundamental mechanisms of immune regulation that have implications far beyond cancer. The understanding of these inhibitory pathways has informed the development of immunotherapies for autoimmune diseases, infectious diseases, and transplant rejection, where the goal is to suppress rather than enhance immune responses.

The success of checkpoint inhibitors has also stimulated research into additional immune checkpoint molecules and inhibitory pathways that could be targeted for therapeutic purposes. LAG-3 (lymphocyte-activation gene 3), TIM-3 (T cell immunoglobulin and mucin domain-3), and TIGIT (T cell immunoreceptor with immunoglobulin and ITIM domains) are among the inhibitory receptors currently being explored as potential checkpoint inhibitor targets.

The field of cancer immunotherapy continues to evolve rapidly, with new combinations, new targets, and new approaches being developed to improve patient outcomes. The legacy of Allison and Honjo's work provides the foundation for these ongoing advances and will continue to shape the future of cancer treatment for generations to come.


\section{Modern Developments}

Allison and Honjo's checkpoint inhibitor work has led to modern cancer immunotherapy. Checkpoint inhibitors (pembrolizumab, nivolumab, ipilimumab) have transformed treatment of melanoma, lung cancer, and many other cancers. Understanding of immune resistance mechanisms has led to combination immunotherapies. Understanding of biomarkers (PD-L1, tumor mutational burden) has enabled patient selection. Understanding of irAEs (immune-related adverse events) has improved management of side effects.


\section{Bibliography}

\begin{thebibliography}{9}
\bibitem{nobel-2018}
Nobel Prize Outreach, ``The Nobel Prize in Physiology or Medicine 2018,''\emph{NobelPrize.org}.\\
\url{https://www.nobelprize.org/prizes/medicine/2018/summary/}

\bibitem{lecture-2018}
James P. Allison, ``Nobel Lecture,''\emph{NobelPrize.org}. Winner-authored primary source; downloadable from the official Nobel Prize archive.\\
\url{https://www.nobelprize.org/prizes/medicine/2018/allison/lecture/}
\end{thebibliography}
